\documentclass[11pt,a4paper]{article}

%----------------------------------------------------------------------------------------
%   PAQUETES Y CONFIGURACIÓN GENERAL
%----------------------------------------------------------------------------------------

% --- Configuración de idioma y codificación ---
\usepackage[utf8]{inputenc}
\usepackage[T1]{fontenc}
\usepackage[spanish]{babel}
\usepackage{csquotes} % Para comillas tipográficas correctas con biblatex

% --- Tipografía y diseño ---
\usepackage{lmodern} % Usa la fuente Latin Modern
\usepackage{microtype} % Mejoras tipográficas sutiles (espaciado, justificación)
\usepackage[a4paper, margin=1cm]{geometry} % Márgenes más equilibrados
\usepackage{setspace} % Permite ajustar el interlineado, ej. \onehalfspacing
\usepackage{parskip} % Párrafos separados por espacio vertical en lugar de sangría

% --- Colores y gráficos ---
\usepackage[dvipsnames]{xcolor} % Permite definir y usar colores
\usepackage{graphicx}

% --- Encabezados y pies de página ---
\usepackage{tabularx}
\usepackage{fancyhdr}
\pagestyle{fancy}
\fancyhf{} % Limpia todos los campos de encabezado y pie de página
\renewcommand{\headrulewidth}{0pt} % Sin línea en el encabezado
\renewcommand{\footrulewidth}{0pt} % Sin línea en el encabezado

% --- Títulos de sección ---
\usepackage{titlesec}
\definecolor{AccentColor}{HTML}{2A52BE} % Un azul profesional (e.g., Cerulean Blue)
\titleformat{\section}
  {\Large\bfseries\color{AccentColor}}
  {\thesection}
  {0.1cm}
  {}
  [\color{AccentColor}\titlerule]
\titleformat{\subsection}
  {\large\bfseries}
  {\thesubsection}
  {0.1cm}
  {}

\titlespacing*{\section}{0pt}{0.1em}{0.2em}
\titlespacing*{\subsection}{0pt}{0.6em}{0.3em}

% --- Listas ---
\usepackage{enumitem}
\usepackage{multicol}

% --- Hipervínculos y referencias ---
\usepackage{hyperref}
\hypersetup{
    colorlinks=true,
    linkcolor=AccentColor,
    urlcolor=AccentColor,
    citecolor=AccentColor,
    pdftitle={Taxonomía de Problemas de Evasión en Grafos},
    pdfauthor={Dante Culaciati}
}
\usepackage[style=numeric, backend=biber, sorting=none]{biblatex}
\renewcommand*{\bibfont}{\footnotesize}
\addbibresource{referencias.bib}

% --- Bookmarks PDF ---
\usepackage{bookmark}

%----------------------------------------------------------------------------------------
%   INFORMACIÓN DEL DOCUMENTO
%----------------------------------------------------------------------------------------

\title{
    \vspace{-1.5cm} % Ajuste para subir el título
    \hrule
    \vspace{0.05cm}
    \textbf{Taxonomía de Problemas de Evasión en Grafos}
    \vspace{0.2cm}
    \hrule
    \vspace{-0.5cm}
}

\author{
    \begin{tabular}{lll}
        \textbf{Alumno:} Dante Culaciati & \textbf{Mentor:} Tomás Delvechio & \textbf{Director:} Esteban Feuerstein \\[0.1cm]
        \multicolumn{3}{l}{\textbf{Lugar de trabajo:} Grupo Algoritmos, Instituto de Ciencias de la Computación,}\\
        \multicolumn{3}{l}{Facultad de Ciencias Exactas y Naturales, Universidad de Buenos Aires}
    \end{tabular}
}

\date{\vspace{-0.2cm} \today} % Usa la fecha de compilación actual

%----------------------------------------------------------------------------------------
%   DOCUMENTO
%----------------------------------------------------------------------------------------

\pagenumbering{gobble}

\begin{document}

\maketitle
\vspace{-1.2cm}

\section*{Introducción}
Los algoritmos de \textit{community detection} permiten identificar  \textit{clusters} o comunidades, o sea grupos de nodos en un grafo relacionados entre sí  a través de muchas aristas, y comparativamente pocas aristas a nodos de otras comunidades \cite{Fortunato-2010}. 
Entre muchas otras aplicaciones, el concepto es utilizado en plataformas de uso masivo, asociado a la identificación de patrones de información de los usuarios, potencialmente vulnerando sus derechos o deseos de privacidad y confidencialidad.

Para mitigar parcial o completamente estos efectos, aparece la idea de \textit{evasión de comunidades}.
La idea general es intentar modificar, engañar, saturar o evadir los resultados del algoritmo de community detection (CD) por medio de la alteración estratégica de las propiedades estructurales del grafo de entrada. Es decir, agregar, eliminar o modificar arcos o atributos de los elementos de la red. 
El nombre ``evasión'' revela de por sí un cierto sesgo, pues veremos que muchos otros problemas pueden ser vistos desde esta misma perspectiva,
en particular \textbf{\emph{\citetitle{bernini-2024}}} planteado por \textcite{bernini-2024} y principal motivación para nuestra propuesta.

Trabajos previos resuelven el problema de \textit{hiding} o \textit{community deception} (planteado originalmente por \textcite{8118127}) a través de algoritmos genéticos \cite{LIU2022123}, heurísticas \cite{Waniek_2018},
algoritmos \textit{greedy} \cite{9377481}, métodos de optimización por gradiente \cite{silvestri2025communitymembershiphidinggradientbased}, autocodificadores de grafos \cite{LIU2022109495}, entre otros. \textcite{Su_2024} exploran ampliamente sobre los conceptos de \textit{detection} y \textit{deception}.
En el paper de Bernini, Silvestri y Tolomei, se formula el problema de \textit{Community Membership Hiding} como un objetivo contrafactual
restringido en grafos, se plantea un proceso de decisión de Márkov y se lo resuelve con deep reinforcement learning (basado en \cite{Chen_2022}).

En ese trabajo, se imponen algunas ``restricciones'', como ser la elección de un \textbf{único} nodo objetivo $u$ que será ocultado, la asignación de una  \textbf{única} comunidad a cada nodo (\textit{non-overlapping communities}), el algoritmo $f(\cdot)$ de CD es considerado una ``\textbf{caja negra}'', entre otras.

\section*{Objetivos}
% (muy abstracto? falta definir CLARAMENTE lo que vamos a hacer?)

El objetivo principal de este proyecto es expandir el espacio de problemas de evasión y analizar escenarios alternativos en el marco de \textit{community evasion}, partiendo de
la base determinada por \textcite{bernini-2024}.
En particular, nos proponemos analizar diversas variantes de la idea. Algunas de ellas son las siguientes (no necesariamente serán analizadas todas durante el trabajo de esta beca):

\begin{multicols}{2}
\begin{itemize}[leftmargin=*, label=\textbullet, itemsep=0em, topsep=0pt]
    \item Inspirado por \cite{8118127}, extendemos el método de \textit{community membership hiding} y lo corremos para cada nodo en la comunidad objetivo.
    \item Trabajamos con \textit{overlapping communities} \cite{Xie_2013}.
    \item Consideramos $f(\cdot)$ como ``white-box'' \cite{jin2021surveycommunitydetectionapproaches}.
    \item Implementamos \textit{node attributes} \cite{Yang_2013}.
    \item Trabajamos con información parcial sobre el grafo.
    \item Consideramos distintos conjuntos de acciones posibles.
    \item Cambiamos el ``scope'' de \textit{community membership hiding} (por ejemplo, buscamos que $u$ no sólo forme parte de otra comunidad, sino que no comparta comunidad con un conjunto selecto de nodos).
    \item Probamos diferentes métodos de optimización \cite{lucic2022, trappolini2022sparseviciousattacksgraph}.
    \item Implementamos \textit{multi-agent reinforcement learning} (MARL) \cite{4445757}.
    \item Variamos las métricas utilizadas para medir distancia entre entrada y salida a $f(\cdot)$.
\end{itemize}
\end{multicols}

\section*{Factibilidad}
% (Charlar bien con Esteban y el mentor. Además, dónde o cómo lo meto a Gabriel? Como co-director?)

Esteban Feuerstein, propuesto como director, es un reconocido investigador y profesor en el área de Diseño y Análisis de Algoritmos
particularmente en el área de grafos y redes, y en la actualidad especializándose en aplicaciones a redes sociales. El Esp. Tomás Delvechio, propuesto como mentor, es integrante de un grupo de investigación de la Universidad Nacional de Luján, trabajando en problemas de búsqueda a gran escala en la intersección de la recuperación de información, las redes sociales y la ciencia de datos.
Cabe destacar que del trabajo a encarar será partícipe también el Dr. Gabriel Tolosa, profesor asociado e investigador en Ciencias de la Computación en la Universidad Nacional de Luján, y supervisor del grupo de investigación del Esp. Tomás Delvechio.

Disponemos del código del paper original (\url{https://github.com/AndreaBe99/community_membership_hiding}) que incluye, entre otros recursos,
varios datasets públicamente disponibles sobre los cuales testear. Además, tenemos contacto directo con los autores originales del paper,
que pueden darnos las herramientas necesarias para resolver ciertos problemas si surgiese la necesidad.

\section*{Cronograma}
% (muy ambiguo?? Repasar. Definir cuales serían los conceptos simples y complejos. Está mal listar MARL explícitamente? Quizás es muy complicado y me termina jugando en contra)
\begin{itemize}[leftmargin=*, label=\textbullet, itemsep=0.5ex]
    \item \textbf{1\textsuperscript{er} trimestre:} Revisión de la literatura, replicación de estudios previos.
    \item \textbf{2\textsuperscript{do} trimestre:} Experimentación inicial. Implementación de conceptos ``simples'' (straightforward).
    \item \textbf{3\textsuperscript{er} trimestre:} Implementación de conceptos complejos, a partir de la experiencia obtenida en el primer semestre.
    \item \textbf{4\textsuperscript{to} trimestre:} Revisión y prueba sobre nuevos datos. Elaboración de conclusiones y redacción del informe final.
\end{itemize}

% (SIN LOS COMENTARIOS LLEGO A DOS CARILLAS)

\begingroup
\tiny
\nocite{*} % Incluye todas las referencias del .bib, citadas o no
\printbibliography[title={Referencias}]
\endgroup

\end{document}